\section{Demo Scenario}

\subsection{Scenario 1: Admin cấu hình registry}

Admin vào \texttt{/login}, chọn role Admin, kết nối ví bootstrap admin và chuyển sang Sepolia. Nếu chưa có contract, Admin deploy \texttt{CredentialRegistry} bằng Remix, copy địa chỉ contract về portal và ký thông điệp cấu hình. Sau khi lưu contract, Admin nhập ví University và gọi \texttt{addIssuer}.

\subsection{Scenario 2: University cấp credential}

University vào \texttt{/login}, chọn role University và kết nối ví đã được authorize. Portal kiểm tra \texttt{authorizedIssuers[wallet]}. University nhập thông tin sinh viên, degree, field, graduation year và danh sách học phần. Khi bấm \textit{Sign, anchor \& issue credential}, app tạo Merkle root, tính UID, yêu cầu ví ký UID và gửi transaction anchor lên Sepolia.

\subsection{Scenario 3: Student tạo disclosure}

Student vào \texttt{/student}, kết nối ví hoặc nhập địa chỉ ví để xem credential. Student chọn credential, tick các học phần muốn chia sẻ, rồi bấm \textit{Generate Merkle proofs}. App tạo presentation JSON chỉ chứa các học phần được chọn, proof tương ứng, metadata và chữ ký issuer. Student có thể copy JSON hoặc tạo share code.

\subsection{Scenario 4: Verifier xác minh}

Verifier vào \texttt{/verify}, nhập share code hoặc paste JSON. Hệ thống kiểm tra UID, chữ ký, Merkle proofs, issuer authorization, anchor và revocation. Kết quả cuối cùng là \textit{Credential verified} nếu tất cả check pass; nếu một trường bị sửa hoặc credential bị revoke, giao diện hiển thị kiểm tra nào thất bại.
